\chapter{Spatial Econometrics}

\paragraph*{Experimental feature.} Every \texttt{spatial\_*} command in this
chapter requires Hayashi to be built with \texttt{--features experimental}.

\section{Why spatial?}

When observations are located in space or connected by a network,
disturbances or outcomes can spill over from one unit to its
neighbours. Ignoring this dependence produces biased standard errors,
coefficient bias, and misleading inference.

\section{The spatial weights matrix}

Every spatial model in \hayashi{} requires a row-standardized weights
matrix $W$. If the data are sorted by unit, $W_{ij}$ is the weight that
unit $j$ receives when predicting the outcome of unit $i$.

\begin{lstlisting}[caption={Spatial weights matrix}]
// 3 regions with mutual neighbours
let W = [[0, 0.5, 0.5],
         [0.5, 0, 0.5],
         [0.5, 0.5, 0]]
\end{lstlisting}

For larger grids, generate $W$ from the DataFrame or load it from a file.

\section{Spatial SAR}

The Spatial Autoregressive model is

\[
y = \rho W y + X\beta + \varepsilon
\]

where $\rho$ is the spatial autoregressive parameter. Estimation uses
concentrated maximum likelihood: for each value in a grid of $\rho$,
transform the dependent variable and run OLS; then select the $\rho$
that maximises the concentrated log-likelihood.

\begin{lstlisting}[caption={Spatial SAR}]
let n = 100
let W = makeweights(df, "region", type="rook")  // or list of lists
generate df x1 = rnormal(n)
generate df y = 0.5*x1 + rnormal(n)
generate df wy = spatial_lag(W, y)
generate df y = y + 0.3*wy

let m_sar = spatial_sar(y ~ x1, df, w=W)
print(m_sar)
\end{lstlisting}

The output reports $\rho$, $\beta$, their standard errors, and a
likelihood-ratio test for $\rho = 0$.

\section{Spatial SEM}

The Spatial Error Model is

\[
y = X\beta + u, \qquad u = \lambda W u + \varepsilon
\]

Errors are spatially correlated, but the outcome itself is not directly
influenced by neighbours' outcomes.

\begin{lstlisting}[caption={Spatial SEM}]
let m_sem = spatial_sem(y ~ x1, df, w=W)
print(m_sem)
\end{lstlisting}

\section{Spatial Durbin and panel extensions}

The Spatial Durbin model adds spatially lagged regressors:

\[
y = \rho W y + X\beta + W X \theta + \varepsilon
\]

It nests SAR ($\theta = 0$) and the spatial lag model ($\rho = 0$).

\begin{lstlisting}[caption={Spatial Durbin}]
let m_sdm = spatial_durbin(y ~ x1, df, w=W)
print(m_sdm)
\end{lstlisting}

For panel data with fixed effects, use the panel variants:

\begin{lstlisting}[caption={Spatial panel SAR}]
let m_panelsar = spatial_panel_sar(y ~ x1, df, w=W, id="region")
print(m_panelsar)
\end{lstlisting}

\section{Interpreting direct and indirect effects}

In a Spatial Durbin model the coefficient on $x$ is not the total effect.
The total effect on unit $i$ includes the direct impact $\beta$ plus the
indirect impact that travels through the network via $\theta$ and $\rho$.

\hayashi{} reports direct, indirect, and total effects in the \texttt{m\_sdm}
summary when available.

\section{Practical notes}

\begin{itemize}
\item Always check that $W$ is row-standardised; non-standardised $W$
      distorts the interpretation of $\rho$ and $\lambda$.
\item Larger samples are needed for spatial MLE than for OLS.
\item Cross-section spatial models assume a stable network. Panel spatial
      models add fixed effects but still require a time-invariant $W$.
\item The \texttt{spatial\_durbin\_error} model combines spatially lagged
      regressors with a spatial error term.
\end{itemize}
